% FORMALIZACION_REUNIDA: CG-015/P09 del inventario REV05.
% Propietarios: sections/nucleo.tex (Cayley APP), libro/toro27.tex:123--256,
% PAQUETE_ARTICULO_I_INTEGRACION_ESTADO_CAMPOS_20260919/deltas/integracion/
% APPRouteWinding.lean (desplazamiento y cociclo entero).
% Esta insercion explicita el cambio de soporte y las fibras aritmeticas.
\section{Refinamiento bidimensional de la APP y cociclo de acarreo}
\label{sec:refinamiento-app-27-9}

La realización de la APP sobre un toro de nueve posiciones por coordenada
y su ampliación a veintisiete posiciones se relacionan mediante división
euclídea. La ampliación conserva dos direcciones independientes y asigna
un dígito ternario adicional a cada coordenada. Por ello cada posición
de la APP inicial posee nueve levantamientos en el soporte ampliado.
Los dos dígitos adicionales conservan la posición dentro de la fibra,
mientras que la dirección del transporte y el régimen del lector TRIT
permanecen como coordenadas propias del estado enriquecido.

Para $m\geq1$, escribimos
\[
 T_m=(\mathbb Z/m\mathbb Z)^2.
\]
En los cálculos de esta sección se eligen representantes
$0,\ldots,m-1$. La presentación nonádica $1,\ldots,9$ imprime la
clase cero mediante el signo $9$. Esta convención permite leer la tabla
APP y efectuar las divisiones enteras con una regla única de residuos.

\subsection{Proyección, fibras y coordenadas ternarias de refinamiento}

La proyección y la sección de representantes son
\[
 p:T_{27}\longrightarrow T_9,\qquad
 p(x,y)=(x\bmod9,y\bmod9),
\]
\[
 s:T_9\longrightarrow T_{27},\qquad s(r,t)=(r,t),
 \quad 0\leq r,t<9.
\]
La proyección es un homomorfismo sobreyectivo de grupos aditivos. Su
núcleo es
\[
 \ker p=\{(9a,9b):a,b\in\mathbb F_3\}
        \simeq\mathbb F_3^2.
\]
La aplicación de coordenadas
\begin{equation}
 \Psi:T_9\times\mathbb F_3^2\longrightarrow T_{27},\qquad
 \Psi((r,t),(a,b))=(r+9a,t+9b)
 \label{eq:app-refinamiento-coordenadas}
\end{equation}
utiliza representantes $0\leq a,b<3$.

\begin{proposition}[Descomposición única de las posiciones ampliadas]
\label{prop:app-refinamiento-fibras}
La aplicación $\Psi$ es biyectiva. Cada fibra de $p$ contiene nueve
posiciones y la totalidad del soporte satisface
\[
 |T_{27}|=27^2=9^2\cdot3^2=729.
\]
Las dos coordenadas de $\mathbb F_3^2$ son los dos trits posicionales
que describen el refinamiento de las dos coordenadas de la APP.
\end{proposition}
\begin{proof}
Para $0\leq x,y<27$, las divisiones euclídeas
\[
 x=r+9a,\qquad y=t+9b,
 \quad 0\leq r,t<9,\quad0\leq a,b<3
\]
existen y son únicas. Estas fórmulas construyen la inversa de $\Psi$.
Fijados $r,t$, los nueve pares $(a,b)$ son distintos y agotan la
fibra. El censo resulta al multiplicar las $81$ posiciones residuales
por sus nueve levantamientos.
\end{proof}

Por ejemplo, la posición $(5,5)$ posee la fibra
\[
 \{(5+9a,5+9b):0\leq a,b<3\}.
\]
La fila de esta fibra con $b=1$ es
$(5,14),(14,14),(23,14)$. Cada posición conserva el mismo par residual
$(5,5)$ y un par distinto de cocientes. La multiplicidad nueve expresa
dos elecciones ternarias independientes. La lectura topológica y
temporal del TRIT conserva, además, su propio dominio operatorio.

\subsection{Composición aditiva y memoria del cruce modular}

Para $r,t\in\{0,\ldots,8\}$, definimos
\[
 r\oplus t=(r+t)\bmod9,\qquad
 c_9(r,t)=\left\lfloor\frac{r+t}{9}\right\rfloor.
\]
Transportada por $\Psi$, la adición en una coordenada de
$\mathbb Z/27\mathbb Z$ es
\begin{equation}
 (r,a)\boxplus(t,b)
   =\bigl(r\oplus t,\ a+b+c_9(r,t)\pmod3\bigr).
 \label{eq:app-refinamiento-adicion}
\end{equation}
En dos coordenadas, la misma fórmula se aplica por separado a cada eje.

\begin{proposition}[Cociclo de la extensión posicional]
\label{prop:app-refinamiento-cociclo}
El acarreo satisface
\begin{equation}
 c_9(r,t)+c_9(r\oplus t,u)
 =c_9(t,u)+c_9(r,t\oplus u)
 =\left\lfloor\frac{r+t+u}{9}\right\rfloor.
 \label{eq:app-refinamiento-cociclo}
\end{equation}
Por tanto, \eqref{eq:app-refinamiento-adicion} es asociativa y
reproduce exactamente el grupo aditivo del soporte ampliado.
\end{proposition}
\begin{proof}
La igualdad $r+t=(r\oplus t)+9c_9(r,t)$ da
\[
 \left\lfloor\frac{r+t+u}{9}\right\rfloor
 =c_9(r,t)+\left\lfloor\frac{(r\oplus t)+u}{9}\right\rfloor.
\]
La agrupación $r+(t+u)$ proporciona el otro miembro. Sustituir
$x=r+9a$ e $y=t+9b$ en $x+y$ prueba
\eqref{eq:app-refinamiento-adicion}; la reducción de su segundo
componente módulo tres corresponde a reducir $x+y$ módulo veintisiete.
La identidad del cociclo prueba directamente la asociatividad de esa
escritura por residuos y cocientes.
\end{proof}

En particular, $8+1=0+9\cdot1$ muestra el cruce elemental: el residuo
retorna a la clase cero, impresa como $9$ en la tabla nonádica, y el
primer cociente aumenta en una unidad. Si se conserva únicamente el
residuo, se pierde el dato que distingue $0$, $9$ y $18$ en el soporte
ampliado; el estado levantado distingue esas tres posiciones.

La extensión
\[
 0\longrightarrow\mathbb Z/3\mathbb Z
 \xrightarrow{\ a\mapsto9a\ }
 \mathbb Z/27\mathbb Z
 \longrightarrow\mathbb Z/9\mathbb Z\longrightarrow0
\]
es no escindida como extensión de grupos. En efecto, cualquier
levantamiento de la clase $1$ es $1+9a$, de orden veintisiete. Una
sección homomorfa tendría que enviar un elemento de orden nueve a uno
cuyo orden dividiera nueve. La sección de representantes posee, en
cambio, el defecto exacto $9c_9$, que la fórmula anterior conserva.

\subsection{Compatibilidad de la hoja multiplicativa}

La misma descomposición conserva la segunda operación de la APP. Para
una coordenada, definimos
\[
 r\odot t=rt\bmod9,\qquad
 m_9(r,t)=\left\lfloor\frac{rt}{9}\right\rfloor.
\]
El producto levantado es
\begin{equation}
 (r,a)\boxtimes(t,b)
 =\bigl(r\odot t,\ m_9(r,t)+rb+ta\pmod3\bigr).
 \label{eq:app-refinamiento-producto}
\end{equation}
\begin{proof}
La expansión entera proporciona
\[
 (r+9a)(t+9b)=rt+9(rb+ta)+81ab.
\]
El último término se anula módulo veintisiete. Al sustituir
$rt=(r\odot t)+9m_9(r,t)$ se obtiene la fórmula. La unicidad de la
división euclídea prueba que sus dos componentes recuperan el producto
residual y su cociente de refinamiento.
\end{proof}
Por ejemplo, $8\cdot8=64\equiv10\pmod{27}$ se representa mediante
$(1,1)$: el residuo es $1$ y
$m_9(8,8)=7\equiv1\pmod3$. Así, suma y producto conservan sus
respectivos acarreos dentro de un mismo soporte ampliado.

\subsection{Direcciones y realización del transporte}

Sea $\mathcal D=\{N,S,E,O\}$ el registro de direcciones y
$\delta_d\in\mathbb Z^2$ el desplazamiento unitario asociado a $d$.
En el soporte con dirección definimos
\[
 S_m(z,d)=(z+\delta_d\pmod m,d).
\]
La proyección conserva explícitamente la dirección,
$p_{\mathcal D}(z,d)=(p(z),d)$, y cumple
\begin{equation}
 p_{\mathcal D}S_{27}=S_9p_{\mathcal D}.
 \label{eq:app-refinamiento-direccion}
\end{equation}
El cálculo del cociente usa, para un desplazamiento entero $v$,
\[
 c_9(r;v)=\left\lfloor\frac{r+v}{9}\right\rfloor,
 \qquad r'=r+v-9c_9(r;v),\qquad
 a'=a+c_9(r;v)\pmod3.
\]
La fórmula incluye los desplazamientos negativos. Por ejemplo,
$r=0,v=-1$ da $r'=8,c_9=-1$; el retorno inverso restituye el cociente.

\begin{proposition}[Conjugación del desplazamiento en coordenadas de fibra]
El cambio de coordenadas $\Psi$ transforma $S_{27}$ en el
desplazamiento residual $S_9$ acompañado por la actualización de los
dos cocientes que acaba de definirse. Esta transformación es biyectiva
y conserva la dirección. Sus potencias conservan la misma identidad
para cualquier longitud finita del transporte.
\end{proposition}
\begin{proof}
En cada eje se sustituye $x=r+9a$ en $x+v$ y se realiza la división
euclídea. La dirección permanece fija en la definición de $S_m$.
La inversa utiliza $-v$ y el cociclo, por lo que recupera el estado
inicial. Componer la identidad una vez por cada transición prueba
el enunciado para todas las potencias.
\end{proof}

La realización de Hilbert conserva esta descomposición. Sobre las bases
de posiciones, el operador
\[
 \mathcal U\delta_{\Psi(z,a),d}
       =\delta_{z,d}\otimes\delta_a
\]
define una isometría sobreyectiva
\[
 \ell^2(T_{27}\times\mathcal D)
 \simeq\ell^2(T_9\times\mathcal D)\otimes\ell^2(\mathbb F_3^2).
\]
El desplazamiento conjugado es el operador de permutación con acarreo
descrito en la proposición. Un factor interno adicional $E$ permanece
intacto al tensorizar con $I_E$. Si un lector de fase depende sólo de
la proyección y un operador direccional actúa con la misma matriz en
todas las posiciones, sus cuadrados de compatibilidad también
conmutan. Un lector que utiliza el cociente se realiza, precisamente,
en el factor de fibra conservado. Estas condiciones explicitan la
relación entre el soporte ampliado y los factores del operador sobre
el toro, con las acciones internas declaradas en su construcción.

\subsection{Seis dígitos ternarios y estructuras algebraicas distintas}

Cada coordenada de $T_{27}$ posee tres dígitos de base tres. Su escritura
conjunta define una biyección de conjuntos
\[
 T_{27}\longleftrightarrow\mathbb F_3^6.
\]
Esta biyección publica seis dígitos posicionales. La adición del toro
contiene acarreos entre los dígitos y la adición vectorial de
$\mathbb F_3^6$ se efectúa componente a componente. Por ejemplo, la
suma posicional $2+1=3$ transforma $(2,0,0)$ en $(0,1,0)$, mientras
que la suma vectorial ternaria daría $(0,0,0)$. El toro contiene
elementos de orden veintisiete; el espacio vectorial ternario tiene
exponente tres. Los dos objetos comparten una codificación finita,
con operaciones tipadas diferentes. La fibra de palabras ternarias,
el soporte de posiciones y las rectas de una superficie cúbica
conservan, de este modo, sus construcciones y sus mapas específicos.
